\documentclass[11pt,letterpaper]{article}

%====================================================================
% 1. MATEMÁTICA Y FUENTES
%====================================================================

% Compilador: XeLaTeX
\usepackage{fontspec}
\usepackage{amsmath}
\usepackage{mathtools}
\usepackage{amsthm}
\usepackage{amssymb}
\usepackage{unicode-math}

%====================================================================
% 2. IDIOMA
%====================================================================

\usepackage{polyglossia}
\setmainlanguage{spanish}

\usepackage{csquotes}
\usepackage{etoolbox}

%====================================================================
% 3. DISEÑO Y FORMATO GENERAL
%====================================================================

\usepackage[margin=1in]{geometry}

\usepackage{microtype}
\usepackage{setspace}
\usepackage{enumitem}
\usepackage{ragged2e}

% Interlineado
\setstretch{1.1}

% Configuración de párrafos
\setlength{\parindent}{1.2em}
\setlength{\parskip}{0pt}
\setlength{\emergencystretch}{2em}

% Configuración de listas
\setlist{
    itemsep=2pt,
    topsep=4pt
}

%====================================================================
% 4. TABLAS
%====================================================================

\usepackage[table]{xcolor}

\usepackage{booktabs}
\usepackage{tabularx}
\usepackage{longtable}
\usepackage{array}

%====================================================================
% 5. FIGURAS Y ELEMENTOS GRÁFICOS
%====================================================================

\usepackage{graphicx}
\usepackage{float}
\usepackage{caption}
\usepackage{pdflscape}

% Configuración de títulos de tablas y figuras
\captionsetup{
    font=small,
    labelfont=bf,
    labelsep=period,
    skip=6pt
}

%====================================================================
% 6. CÓDIGO Y LISTADOS
%====================================================================

\usepackage{listings}

\lstset{
    language=Python,
    basicstyle=\ttfamily\small,
    frame=lines,
    keywordstyle=\bfseries,
    commentstyle=\itshape\color{gray},
    showstringspaces=false,
    columns=fullflexible,
    xleftmargin=1em
}

%====================================================================
% 7. ENCABEZADOS Y PIES DE PÁGINA
%====================================================================

\usepackage{fancyhdr}

\setlength{\headheight}{15pt}

\pagestyle{fancy}
\fancyhf{}

\renewcommand{\headrulewidth}{0pt}

\fancyhead[R]{
    \small\itshape Informe de avance del proyecto
}

\fancyfoot[C]{
    \small\thepage
}

%====================================================================
% 8. BIBLIOGRAFÍA
%====================================================================

% Procesamiento mediante Biber
\usepackage[
    backend=biber,
    style=authoryear,
    sorting=nyt
]{biblatex}

% Archivo de referencias bibliográficas
\addbibresource{references.bib}

%====================================================================
% 9. ENLACES Y REFERENCIAS CRUZADAS
%====================================================================

\usepackage{xurl}

\usepackage[
    colorlinks=true,
    linkcolor=blue!40!black,
    citecolor=blue!40!black,
    urlcolor=blue!40!black,
    pdfusetitle
]{hyperref}

\hypersetup{
    pdftitle={Informe de avance del proyecto},
    pdfauthor={Anthonny Flores Rojas,
               Andrey Gonzalez Bastos,
               Diego Vega}
}

\usepackage[capitalise]{cleveref}

% Nombres utilizados en las referencias cruzadas
\crefname{theorem}{Teorema}{Teoremas}
\crefname{proposition}{Proposición}{Proposiciones}
\crefname{lemma}{Lema}{Lemas}
\crefname{corollary}{Corolario}{Corolarios}
\crefname{definition}{Definición}{Definiciones}
\crefname{example}{Ejemplo}{Ejemplos}
\crefname{conjecture}{Conjetura}{Conjeturas}
\crefname{section}{Sección}{Secciones}
\crefname{table}{Tabla}{Tablas}
\crefname{figure}{Figura}{Figuras}
\crefname{equation}{ecuación}{ecuaciones}

%====================================================================
% 10. TEOREMAS Y ENTORNOS MATEMÁTICOS
%====================================================================

% Numeración de teoremas por sección
\theoremstyle{plain}

\newtheorem{theorem}{Teorema}[section]
\newtheorem{proposition}[theorem]{Proposición}
\newtheorem{lemma}[theorem]{Lema}
\newtheorem{corollary}[theorem]{Corolario}
\newtheorem{conjecture}[theorem]{Conjetura}

\theoremstyle{definition}

\newtheorem{definition}[theorem]{Definición}
\newtheorem{example}[theorem]{Ejemplo}

\theoremstyle{remark}

\newtheorem{remark}[theorem]{Observación}

% Numeración de ecuaciones por sección
\numberwithin{equation}{section}

% Nombre de las demostraciones
\gappto\captionsspanish{%
    \renewcommand{\proofname}{Demostración}
}

%====================================================================
% 11. COMANDOS MATEMÁTICOS
%====================================================================

\newcommand{\R}{\mathbb{R}}
\newcommand{\Z}{\mathbb{Z}}
\newcommand{\E}{\mathbb{E}}

\renewcommand{\P}{\mathbb{P}}

\newcommand{\dd}{\,\mathrm{d}}

\newcommand{\norm}[1]{%
    \lVert #1\rVert
}

\newcommand{\Sh}{\mathcal{S}}

\DeclareMathOperator{\LSE}{LSE}

%====================================================================
% 12. ENTORNOS MATEMÁTICOS EN ESPAÑOL
%====================================================================

\newtheorem{teorema}[theorem]{Teorema}
\newtheorem{proposicion}[theorem]{Proposición}
\newtheorem{lema}[theorem]{Lema}
\newtheorem{corolario}[theorem]{Corolario}
\newtheorem{definicion}[theorem]{Definición}
\newtheorem{ejemplo}[theorem]{Ejemplo}
\newtheorem{observacion}[theorem]{Observación}

%====================================================================
% 13. COMANDOS PARA MONEDAS
%====================================================================

\newcommand{\eur}{\texttt{EUR}\,}
\newcommand{\gbp}{\texttt{GBP}\,}
\newcommand{\czk}{\texttt{CZK}\,}

%====================================================================
% 14. CLASIFICACIÓN DE DATOS
%====================================================================

% Categoría A: datos publicados
\newcommand{\catA}{\textbf{[A]}}

% Categoría B: cálculos propios
\newcommand{\catB}{\textbf{[B]}}

% Categoría C: datos no disponibles o no verificables
\newcommand{\catC}{\textbf{[C]}}

%====================================================================
% 15. COMANDOS AUXILIARES DEL INFORME
%====================================================================

% Resaltar el estado de un elemento
\newcommand{\estado}[1]{\textbf{#1}}

% Entornos de interlineado sencillo
\newcommand{\sw}{\begin{singlespace}}
\newcommand{\ew}{\end{singlespace}}

% Texto provisional
\newcommand{\ph}[1]{%
    \textcolor{red!70!black}{\textit{[#1]}}
}

% Estados para tablas de avance
\newcommand{\done}{%
    \textcolor{green!45!black}{\textbf{Completado}}
}

\newcommand{\wip}{%
    \textcolor{orange!80!black}{\textbf{En curso}}
}

\newcommand{\todo}{%
    \textcolor{red!70!black}{\textbf{Pendiente}}
}

%====================================================================
% 16. DATOS INSTITUCIONALES
%====================================================================

\newcommand{\Universidad}{Universidad de Costa Rica}
\newcommand{\Escuela}{Escuela de Matemáticas}

\newcommand{\Curso}{
    Valoración de instrumentos financieros (CA0308)
}

\newcommand{\Profesor}{
    Daniel Vargas Carrillo
}

%====================================================================
% 17. DATOS DEL DOCUMENTO
%====================================================================

\title{
    \textbf{Nombre del informe}\\[0.4em]
    \large Informe de anteproyecto
}

\author{%
    {Anthonny Flores Rojas}
    \thanks{
        Escuela de Matemáticas,
        Universidad de Costa Rica.
        \texttt{anthonny.flores@ucr.ac.cr}
    }
    \and
    {Andrey Gonzalez Bastos}
    \thanks{
        Escuela de Matemáticas,
        Universidad de Costa Rica.
        \texttt{andrey.gonzalezbastos@ucr.ac.cr}
    }
    \and
    {Diego Vega}
    \thanks{
        Escuela de Matemáticas,
        Universidad de Costa Rica.
        \texttt{diego.vega@ucr.ac.cr}
    }
}

\date{
    Curso: Valoración de instrumentos financieros (CA0308)\\
    Profesor: Daniel Vargas Carrillo\\[0.3em]
    \today
}

\begin{document}
	
	
	\maketitle
	\tableofcontents   % opcional; muchos papers largos lo incluyen
	
% =====================================================================
\section{Introducción}\label{sec:intro}

Giannozzi y Kirilova \parencite{GIANNOZZI2026107700} investigan por qué una cobertura 
con derivados, útil para fijar precios, puede amplificar el riesgo de una empresa en 
un episodio de estrés. Los contratos cotizados se valoran a mercado a diario y exigen
 margen inicial y de variación en efectivo; si el precio del subyacente se mueve contra 
 la posición, la empresa debe aportar liquidez de inmediato aunque la cobertura sea 
 económicamente razonable \parencite[p.~1]{GIANNOZZI2026107700}. Con datos de empresas 
 energéticas europeas durante la crisis de 2022, los autores documentan que la mayoría 
 financia esas llamadas con deuda, que las empresas sin relaciones previas con prestamistas 
 venden activos, y que las llamadas se asocian con mayor riesgo de incumplimiento, 
 mayor costo de la deuda y menores márgenes \parencite[p.~1]{GIANNOZZI2026107700}. 
 Son resultados de la muestra; en este informe se usan como mecanismo explicativo y 
 no se atribuyen a ČEZ.

La selección de ČEZ, a. s., responde a criterios financieros y no a 
la disponibilidad de datos. Primero, la empresa pertenece a la muestra del artículo: 
figura entre las que informan en qué partida del balance registran los depósitos de margen, 
en este caso «cuentas por cobrar comerciales, margen y colaterales» 
\parencite[Apéndice A3, p.~16]{GIANNOZZI2026107700}. Segundo, genera cerca de 58\% 
de la electricidad checa \parencite[p.~11]{cez_fs_2021} y el Estado posee 69,8\% de su 
capital \parencite[p.~10]{cez_interim_q3_2021}. Tercero, vende por adelantado su generación 
y lo cuantifica antes del corte: 79\%, 49\% y 21\% de la generación de 2022, 2023 y 
2024 al 30 de septiembre de 2021 \parencite[p.~27]{cez_interim_q3_2021}. 
Cuarto, su exposición es sobre todo eléctrica (el gas aporta 5,8\% de su generación 
, \parencite{cez_story_2022}), lo que permite aislar un subyacente. 
Quinto, su balance al 30 de septiembre de 2021 se publicó el 9 de noviembre de 2021, 
antes del corte.

El análisis se limita a lo que podía conocerse el \textbf{23 de febrero de 2022}, 
último día hábil anterior al 24 de febrero, cuando se registran los mayores rendimientos 
diarios del gas y la electricidad \parencite[p.~3]{GIANNOZZI2026107700}. Todos los precios, 
volatilidades y probabilidades usan observaciones hasta esa fecha inclusive; los hechos 
posteriores aparecen solo como evidencia retrospectiva (Sección~\ref{sec:sens}). 

\section{Contexto de la empresa y del mercado en la fecha de corte}

\textbf{Empresa.} ČEZ genera, distribuye y comercializa electricidad y calor, además de gas, servicios energéticos y carbón \parencite[p.~10]{cez_interim_q3_2021}. Las ventas de electricidad fueron CZK 90.496 millones de unos ingresos operativos de CZK 155.955 millones en nueve meses de 2021 (58\% , \parencite[p.~21]{cez_interim_q3_2021}) y el segmento de Generación aportó CZK 24.271 millones de un EBITDA consolidado de CZK 47.520 millones (51\% , \parencite[pp.~24--25]{cez_interim_q3_2021}). Su estrategia combina ventas de electricidad a plazo, compras a plazo de derechos de emisión y cobertura del tipo de cambio de ingresos en euros \parencite{cez_q3_2021}. El costo ya era visible antes del corte: las preventas redujeron en CZK 4.305 millones los ingresos por electricidad de nueve meses porque se liquidaron por debajo del precio de mercado \parencite[p.~21]{cez_interim_q3_2021}. La liquidez de referencia al 30 de septiembre de 2021 era de CZK 8.023 millones en efectivo \parencite[p.~2]{cez_interim_q3_2021} y de CZK 31.400 millones en líneas comprometidas no dispuestas (40.000 menos 8.600 dispuestos; cifra de una diapositiva por confirmar \parencite{cez_q3_2021}), frente a una deuda de CZK 141.245 millones \parencite[pp.~3, 19, 20]{cez_interim_q3_2021}.

\textbf{Mercado.} La Tabla~\ref{tab:mercado} resume los indicadores al corte. El contrato anual alemán para 2023 cotizaba a EUR 156,43/MWh el 23 de febrero de 2022, 82\% más que el 30 de septiembre de 2021 (EUR 85,86) y tres veces el nivel de enero de 2021 (EUR 51,25) . En el mercado spot, el promedio diario de la subasta Day Ahead pasó de EUR 100,66 (CZ) y 96,84 (DE-LU) en 2021 a EUR 183,22 y 178,97 en el cuarto trimestre, y a 166,76 y 147,16 entre el 1 de enero y el 23 de febrero de 2022 . El diferencial medio entre la República Checa y Alemania-Luxemburgo fue de EUR 5,85/MWh (5,7\% del precio medio de Alemania-Luxemburgo), con correlación de 0,976 entre promedios diarios, y los precios coincidieron en 46,6\% de las horas, lo que indica acoplamiento de mercados y un riesgo de base menor que el riesgo de nivel .

\sw
\begin{table}[ht]
\centering
\footnotesize
\begin{tabularx}{\textwidth}{@{}>{\RaggedRight}p{3.0cm}>{\RaggedRight}X>{\RaggedRight}p{3.4cm}@{}}
\toprule
Indicador & Valor y fecha & Fuente y categoría \\
\midrule
Futuro base DE, entrega 2023 & EUR 156,43/MWh (23-feb-2022); 85,86 (30-sep-2021); 51,25 (4-ene-2021) & \textcite{energycharts_futuros_de},  \\
Futuro base DE, entrega 2024 & EUR 118,88/MWh (23-feb-2022) & ídem,  \\
Spot Day Ahead CZ / DE-LU & Último día, 23-feb-2022: EUR 136,63 / 130,09 (promedio diario) & \textcite{energycharts_spot2022},  \\
Gas TTF & EUR 88,85/MWh (23-feb-2022); máximo previo de EUR 182,78 el 21-dic-2021 & \textcite{platts_ttf_2022},  \\
Derechos de emisión EUA & EUR 94,67/t, contrato más cercano (cierre del 23-feb-2022); EUR 61,66 al 30-sep-2021 & \textcite{gasstrategies_2022}; \parencite[p.~18]{cez_interim_q3_2021},  \\
EUR/CZK & CZK 24,860 al 31-dic-2021; cotización del 23-feb-2022 pendiente & \parencite[p.~30]{cez_fs_2021},  \\
Tasas & Repo checa a dos semanas 4,50\% desde el 4-feb-2022 (antes 3,75\%); PRIBOR pendiente & \textcite{cnb_2022},  \\
Regulación & EMIR: compensación central obligatoria sobre EUR 4.000 millones de nocional bruto en materias primas; sin compensación, valoración diaria y garantías & \parencite[p.~18]{GIANNOZZI2026107700}; norma primaria por verificar, / \\
Competencia & ČEZ genera cerca de 58\% de la electricidad checa; competidores y cuotas pendientes & \parencite[p.~11]{cez_fs_2021},  \\
\bottomrule
\end{tabularx}
\caption{Estado del mercado en la fecha de corte.}
\label{tab:mercado}
\end{table}
\ew

El significado financiero de estos niveles es directo. ČEZ vendió 37 TWh de 2022 a EUR 54,1/MWh \parencite{cez_q3_2021}, mientras el mercado para años cercanos ya superaba EUR 150/MWh: las ventas a plazo estaban muy por debajo del precio de mercado, es decir, con pérdida de valoración, antes de que el precio se disparara en febrero. Los derechos de emisión se encarecieron 53\% desde septiembre , y ČEZ estima que cada EUR 1/t se traslada a la electricidad entre EUR 0,6 y 0,7/MWh \parencite{cez_story_2022}. Una repo de 4,50\% encarece además financiar los depósitos de margen. Cuanto mayor es el precio, mayor es la pérdida de la posición vendida y, si hay garantías, mayor el efectivo que debe aportarse, que es el canal que el artículo identifica.




\end{document}